\begin{abstract}
Language agents that work across many sessions keep a running memory: after each session a model rewrites a bounded summary of everything so far.
We ask which claims such a memory loses, and show that it is not random.
We introduce \sys{}, a controlled benchmark in which every claim's content, wording and last mention are held fixed while only the number of times it is mentioned is varied.
Across four domains and four open models, a warning, precondition or objection voiced once is far less likely to survive iterative summarisation than the same claim voiced four times, while a full-context reader shows no such gap; the effect is specific to rewriting, absent under truncation, and survives a schedule that removes the first-mention confound, LLM-naturalised dialogues, and a fixed external reader, so it lives in the memory rather than in the model that reads it.
Prompting the summariser to preserve warnings helps strong models but leaves small models, tight budgets and typed compaction strongly frequency-biased.
We propose \ledger{}, which extracts claims from each new session, consolidates paraphrases in code, and evicts by a type tier rather than by support.
At equal word budget it recovers 6--54 points of accuracy on single-mention critical claims over rolling summarisation and never underperforms the strongest prompt-only baseline on them.
\end{abstract}
